% EUMAS 2026: What Your Agent Topology Is Actually Doing
% Expanded from CAIS v8 (5-6 pages) to 14 pages LNCS
% Draft v1: April 11, 2026
% Double-blind submission — no author names
\documentclass[runningheads]{llncs}

\usepackage{graphicx}
\usepackage{booktabs}
\usepackage{amsmath,amssymb}
% amsthm conflicts with llncs proof environment; omit it
\usepackage[T1]{fontenc}
\usepackage[hyphens]{url}
\usepackage[expansion=false]{microtype}
\usepackage{hyperref}
\usepackage{enumitem}
\usepackage{multirow}
\usepackage{xcolor}

\graphicspath{{figures/}}

% LNCS already defines theorem, lemma, etc.
% We only need custom environments not already provided.
\spnewtheorem{thm}{Theorem}{\bfseries}{\itshape}

\begin{document}

\title{What Your Agent Topology Is Actually Doing:\\
Controlled Experiments on Communication Structure\\
in Multi-Agent Systems}

\titlerunning{What Your Agent Topology Is Actually Doing}

% Double-blind: no authors
\author{Anonymous}
\institute{Anonymous}

\maketitle

\begin{abstract}
Multi-agent AI systems choose communication topologies---star, chain,
ring, fully connected---with no controlled evidence for how topology
affects behavior. We provide five lines of evidence across 1,360 GA
runs and 3,200 LLM API calls. First, we prove that cycle rank and
edge density are algebraically confounded in undirected graphs,
invalidating density-uncontrolled topology comparisons (Theorem~1).
Second, using directed graphs that escape this confound (240~runs), we
show topology independently explains 17\% of diversity variance
($p < 10^{-7}$). Third, an NK dose-response sweep (250~runs) shows
the effect scales $14\times$ from smooth to rugged landscapes. Fourth,
a ring-vs-star comparison at constant cycle rank (160~runs) reveals
that degree distribution contributes at most $\eta^2 \approx 0.13$
within a single $\beta_1$ class---cycle rank matters more than local
connectivity pattern. Fifth, a bridge experiment (320~runs) with
iso-spectral graph families decomposes the topology signal into two
independent mechanisms: cycle rank ($\beta_1$) drives transient
diversity ($\eta^2 = 0.191$ at generation~100, fading by
generation~500), while algebraic connectivity ($\lambda_2$) drives
persistent diversity (cross-family $\eta^2 = 0.265$ at
generation~200, persisting). We validate transfer to real agents:
3,200 GPT-4o-mini API calls show topology modulates LLM solution
diversity with a coupling-dependent sign flip---a single instruction
change reverses which topology produces the most diverse output
($\eta^2 = 0.291$, $p < 10^{-19}$). At least nine deployed systems
chose topologies without the controlled comparisons our confound
theorem demands. Topology is a first-class design variable; current
systems are tuning it blindly.
\end{abstract}

%% ===================================================================
\section{Introduction}
\label{sec:intro}
%% ===================================================================

Every multi-agent system embeds a communication topology. AutoGen uses
round-robin group chat. LangGraph builds DAG pipelines. CrewAI deploys
manager hubs. These are not neutral infrastructure choices---they
determine which agents can see which outputs, shaping the
diversity--convergence tradeoff that governs system behavior. Yet
these choices are made without controlled evidence.

The scale of the problem is becoming clear.
Dochkina~\cite{dochkina2026} ran 25,000~tasks across 256~agents with
eight coordination protocols, finding a 44\% quality spread
($d = 1.86$) between topologies---an effect larger than most prompt
engineering interventions. Moran~\cite{moran2025error} reports
$17\times$ error amplification in unstructured multi-agent
communication. Yang et al.~\cite{yang2025topology} identify topology
as \emph{the} underexplored dimension of multi-agent system design,
calling for controlled experiments. Wang et al.~\cite{wang2026kregular}
show that $K$-regular graphs maximize coordination efficiency, with a
minimum capability threshold $\beta_{\min} \sim 1/K$ suggesting that
connectivity and agent capability trade off.

Topology is doing something---but \emph{what, exactly?} Existing
comparisons conflate multiple graph properties. Comparing a chain
($\beta_1 = 0$) to a fully connected graph ($\beta_1 = \binom{n}{2} -
n + 1$) simultaneously varies cycle rank, density, diameter, and
spectral gap. We cannot know which property drives the observed
effect.

We provide the first controlled answer, bridging from a model system
(genetic algorithms) to real LLM agents. Our contributions:

\begin{enumerate}
  \item \textbf{The Confound Theorem} (Section~\ref{sec:definitions}):
    cycle rank and edge density are algebraically identical for
    undirected graphs ($\beta_1 = |E| - n + 1$), making
    density-uncontrolled topology comparisons fundamentally
    uninterpretable. We show how directed cycle count $\kappa$ escapes
    this confound.

  \item \textbf{Controlled GA experiments} (Sections~\ref{sec:exp1}--\ref{sec:exp4}):
    four experiments totaling 970 runs establish that topology
    independently determines behavioral diversity, the effect scales
    with problem difficulty, cycle rank matters more than degree
    distribution, and two invariants ($\beta_1$, $\lambda_2$) control
    diversity at different timescales.

  \item \textbf{LLM validation} (Section~\ref{sec:exp5}): 3,200
    GPT-4o-mini API calls confirm topology modulates agent diversity,
    with a coupling-dependent sign flip that a single instruction
    change controls.

  \item \textbf{Evidence gap analysis} (Section~\ref{sec:discussion}):
    at least nine recent systems chose topologies without the
    controlled comparisons our confound theorem demands.

  \item \textbf{Invariant hierarchy} (Section~\ref{sec:discussion}):
    we identify three graph invariants---directed cycle count
    $\kappa$, cycle rank $\beta_1$, and algebraic connectivity
    $\lambda_2$---operating at different timescales, providing a
    principled framework for topology selection.
\end{enumerate}

%% ===================================================================
\section{Background: Island Models as Agent Systems}
\label{sec:background}
%% ===================================================================

An island-model genetic algorithm (GA) partitions a population across
sub-populations (\emph{islands}) that evolve independently, exchanging
individuals periodically via \emph{migration}
channels~\cite{tomassini2005spatially}. The mapping to agent
orchestration is direct at the graph level (Table~\ref{tab:mapping}).

\begin{table}[t]
\centering
\caption{Island-model GA concepts and their agent-system equivalents.
The mapping is exact at the graph level (rows 1--5) and analogical at
the node level (row 6). Our object of study---the
diversity--topology relationship---is a graph-level property in both
systems.}
\label{tab:mapping}
\begin{tabular}{@{}lll@{}}
\toprule
\textbf{GA Concept} & \textbf{Agent Equivalent} & \textbf{Level} \\
\midrule
Island & Agent (local state, computation) & Node \\
Migration channel & Message channel & Edge \\
Migration topology & Orchestration graph & Graph \\
Population diversity & Solution diversity & Graph metric \\
Fitness function & Task performance metric & Evaluation \\
\midrule
Population (many individuals) & Single output per agent & Node \\
\bottomrule
\end{tabular}
\end{table}

Both systems are nodes with local state exchanging information along
directed channels. The critical shared structure is the
\emph{topology}---who talks to whom---and its effect on aggregate
diversity. The analogy is weaker at the node level: a GA island
contains a population evolving over many generations, while an agent
produces a single output per round. But the diversity--topology
relationship, our object of study, is a graph-level property in both
cases.

\paragraph{Where the analogy holds and breaks.}
The mapping is strongest for the structural property we study: how
communication topology modulates the diversity of outputs across the
system. GA migration mixes genetic material between islands; agent
message-passing shares solution information between agents. In both
cases, \emph{more} connectivity generally enables \emph{faster}
information spread. The mapping breaks at the mechanism level: GA
migration is blind copying of genotypes, while agent communication
involves semantic understanding and selective incorporation. This
difference is precisely what makes the sign-flip finding
(Section~\ref{sec:exp5}) interesting: the same topology can have
opposite diversity effects depending on the coupling mechanism.

\paragraph{Why GAs as a model system?}
GAs provide cheap, reproducible experimental control. A single
topology comparison requires $\sim$20 seeds to achieve statistical
power; at 500 generations per run, this takes seconds on a laptop.
The equivalent LLM experiment requires thousands of API calls costing
hundreds of dollars. GAs let us sweep parameter spaces (8 topologies
$\times$ 4 difficulty levels $\times$ 20 seeds = 640 runs in
Experiments 1--2) before investing in expensive LLM validation.

\paragraph{Metrics.}
\textbf{GA diversity} is measured as mean pairwise Hamming distance
between all individuals across all islands (normalized by string
length), capturing the fraction of the search space the population
spans. \textbf{GA fitness} is the fitness of the best individual found
across all islands. \textbf{Agent diversity} is mean pairwise cosine
distance on sentence embeddings (all-MiniLM-L6-v2) of agent outputs.
\textbf{Agent fitness} is task performance (test pass rate).

%% ===================================================================
\section{Definitions and the Confound Theorem}
\label{sec:definitions}
%% ===================================================================

We define the graph-theoretic quantities that our experiments control
and measure, together with their agent-system interpretations.

\begin{definition}[Communication topology]
\label{def:topology}
A \emph{communication topology} for $n$~agents is a directed graph
$D = (V, A)$ with $|V| = n$, where each arc $(u, v) \in A$ is a
permitted message channel from agent~$u$ to agent~$v$.
\end{definition}

\noindent
\textit{Agent interpretation:} The topology determines which agents
can see which other agents' outputs. An arc $(u, v)$ means agent~$v$
receives agent~$u$'s output in its context window. No arc means no
information flow.

\begin{definition}[Cycle rank]
\label{def:cyclerank}
The \emph{cycle rank} (first Betti number) of a connected undirected
graph $G = (V, E)$ is
\begin{equation}
  \beta_1(G) = |E| - |V| + 1.
\end{equation}
It counts the number of independent cycles in
$G$~\cite{hatcher2002algebraic}---equivalently, the minimum number of
edges whose removal makes $G$ a tree.
\end{definition}

\noindent
\textit{Agent interpretation:} $\beta_1$ counts independent feedback
loops in the communication structure. A tree ($\beta_1 = 0$) has no
feedback; information flows one way from root to leaves. Each
additional cycle ($\beta_1 = 1, 2, \ldots$) creates a new feedback
pathway where agents can respond to downstream effects of their own
outputs. Capucci and Myers~\cite{capucci2026feedback} formalize this
cost: each cycle requires an independent convergence guarantee
(Lyapunov function) to ensure the system reaches a fixed point.

\begin{definition}[Algebraic connectivity]
\label{def:lambda2}
The \emph{algebraic connectivity} $\lambda_2(G)$ is the second-smallest
eigenvalue of the Laplacian matrix $L = D - A$ of the undirected
graph underlying~$D$, where $D$ is the degree matrix and $A$ the
adjacency matrix. $\lambda_2 > 0$ iff $G$ is
connected~\cite{fiedler1973algebraic}.
\end{definition}

\noindent
\textit{Agent interpretation:} $\lambda_2$ governs the rate at which
information propagates through the communication network. High
$\lambda_2$ means consensus forms quickly (agents rapidly converge to
shared solutions); low $\lambda_2$ means information takes longer to
traverse the network, preserving local diversity for more rounds. A
ring has $\lambda_2 = 2(1 - \cos(2\pi/n))$; a complete graph has
$\lambda_2 = n$.

\begin{definition}[Directed cycle count]
\label{def:kappa}
$\kappa(D)$ counts the number of distinct simple directed cycles in a
directed graph $D$~\cite{johnson1975finding}. Unlike $\beta_1$, this
depends on arc \emph{arrangement}, not just edge count. A DAG has
$\kappa = 0$.
\end{definition}

\noindent
\textit{Agent interpretation:} $\kappa$ counts the number of ways
information can circulate back to its source. Two directed graphs
with the same number of arcs can have vastly different $\kappa$
values (from 0 in a DAG to 47 in our densest cyclic graph), meaning
different amounts of feedback even at identical communication
bandwidth.

\paragraph{Timescale hierarchy.}
\label{rem:hierarchy}
These three invariants operate at different timescales in our
experiments: directed cycle count $\kappa$ determines the
\emph{arrangement} of feedback (Experiment~1), cycle rank $\beta_1$
determines the \emph{amount} of transient diversity
(Experiment~4, generation~100), and algebraic connectivity
$\lambda_2$ determines the \emph{rate} of long-term convergence
(Experiment~4, generation~200+). This suggests a hierarchy:
$\kappa \to \beta_1 \to \lambda_2$, from fast to slow.

We now state the central theoretical result that motivates our
experimental design.

\begin{thm}[Density--Cycle Confound]
\label{thm:confound}
For any connected undirected graph $G$ on $n$ vertices,
$\beta_1(G) = |E| - n + 1$. Therefore, cycle rank and edge density
are not independent variables: every comparison that varies density
simultaneously varies cycle rank, and vice versa.
\end{thm}

\begin{proof}
By elementary algebraic topology~\cite{hatcher2002algebraic}, the
first Betti number of a connected graph equals the number of edges
minus the number of vertices plus the number of connected components.
For a connected graph, this gives $\beta_1 = |E| - |V| + 1 = |E| - n + 1$.
Since edge density $\delta = |E| / \binom{n}{2}$, we have
$\beta_1 = \delta \binom{n}{2} - n + 1$, a linear function of
$\delta$ at fixed $n$. Varying one necessarily varies the other.
\end{proof}

\noindent
\textit{Consequence for agent systems:} Any topology comparison that
varies the number of communication channels (e.g., chain vs.\ fully
connected, star vs.\ mesh) simultaneously varies the number of
feedback loops. Claims about the effect of ``adding more
connections'' cannot distinguish density effects from cycle effects.
To isolate cycle structure, one must either (a)~hold density constant
and vary cycle arrangement (our Experiments~1, 3, 4), or (b)~use
directed graphs where $\kappa$ varies independently of arc count (our
Experiment~1). At least nine recent agent systems fail to make this
distinction (Section~\ref{sec:discussion}).

\paragraph{NK landscapes.}
NK landscapes~\cite{kauffman1993origins} provide tunably rugged
fitness functions over binary strings of length~$N$. Each bit's
fitness contribution depends on $K$~other bits. $K = 0$ gives a
smooth, single-peaked landscape (trivial optimization); increasing $K$
adds local optima, creating a rugged landscape with many competing
strategies. We use $N = 20$ throughout. NK landscapes serve as a
proxy for task complexity: simple tasks (few valid approaches,
unambiguous evaluation) correspond to low $K$; complex tasks (multiple
valid strategies, ambiguous success criteria) correspond to high $K$.

%% ===================================================================
\section{Experiment 1: Directed Cycles at Constant Density}
\label{sec:exp1}
%% ===================================================================

The confound theorem (Theorem~\ref{thm:confound}) shows that
undirected graph comparisons cannot isolate cycle structure from
density. We escape this confound using directed graphs, where the
directed cycle count $\kappa$ varies independently of arc count.

\paragraph{Design.}
Eight directed graphs, all with $n = 8$ vertices and exactly 16~arcs
(constant density $\delta = 2.0$ arcs per vertex), varying only in
directed cycle count:
$\kappa \in \{0, 0, 3, 10, 14, 20, 29, 47\}$. The two $\kappa = 0$
graphs (DAGs) serve as replicated controls. Island-model GA: 50
individuals per island, 500~generations, OneMax fitness function
($K = 0$), 30~seeds per topology, 240~runs total.

% FIGURE: iso-dense-directed-graphs.pdf
% Caption: Eight directed graphs with identical density (16 arcs, 8 nodes)
% but varying directed cycle count kappa from 0 to 47. Despite identical
% bandwidth, cycle structure alone drives the diversity effect.

\paragraph{Results.}
One-way ANOVA: $F(7, 232) = 6.90$, $p = 1.85 \times 10^{-7}$,
$\eta^2 = 0.17$. Topology explains 17\% of diversity variance at
constant density. The correlation between directed cycle count and
final diversity is $r = -0.68$: more directed cycles \emph{accelerate
consensus}. Topologies with more feedback loops cause populations to
converge faster, reducing the system's capacity to maintain
alternative solutions.

\paragraph{Interpretation.}
This result establishes a baseline: even on a trivial landscape
($K = 0$), cycle structure alone---with density perfectly
controlled---produces a medium-to-large effect ($\eta^2 = 0.17$) on
population diversity. The direction is intuitive: more feedback
pathways mean faster information mixing, which accelerates convergence
toward a shared solution. The question is how this effect scales with
problem difficulty.

%% ===================================================================
\section{Experiment 2: NK Dose-Response}
\label{sec:exp2}
%% ===================================================================

Experiment~1 establishes that topology affects diversity on a trivial
problem. Real agent deployments involve complex tasks with multiple
valid approaches. Do topology effects grow or shrink as problems get
harder?

\paragraph{Design.}
NK landscapes with $K \in \{0, 2, 4, 6\}$ (four ruggedness levels),
topologies selected to span the $\kappa$ range from Experiment~1,
10--20 seeds per condition, 250~runs total. All other parameters
(population size, migration rate, number of generations) identical to
Experiment~1.

\paragraph{Results.}
Topology's diversity effect scales monotonically with landscape
ruggedness (Table~\ref{tab:nk-dose}):

\begin{table}[t]
\centering
\caption{Topology effect size ($\eta^2$) on diversity scales
$14\times$ from smooth to rugged NK landscapes. Simple tasks show
negligible topology effects; complex tasks are dominated by topology.}
\label{tab:nk-dose}
\begin{tabular}{@{}ccll@{}}
\toprule
$K$ & $\eta^2$ (diversity) & $p$-value & Interpretation \\
\midrule
0 (smooth) & 0.05 & $> 0.05$ & Topology barely matters \\
2 (moderate) & 0.20 & $< 0.01$ & Moderate effect \\
4 (rugged) & 0.45 & $< 10^{-4}$ & Large effect \\
6 (very rugged) & 0.69 & $< 10^{-6}$ & Topology dominates \\
\bottomrule
\end{tabular}
\end{table}

% FIGURE: eta-squared-vs-k.pdf
% Caption: Topology effect size (eta-squared for diversity) as a function
% of NK landscape ruggedness K. The effect scales 14x from K=0 to K=6,
% demonstrating that complex tasks amplify topology's influence.

\paragraph{Causal chain.}
The mechanism proceeds in stages: topology $\to$ diversity $\to$ mean
fitness, with minimal direct effect on best fitness. At $K = 6$:
topology explains $\eta^2 \approx 0.69$ of diversity variance,
diversity explains $\eta^2 \approx 0.35$ of mean fitness variance,
but topology's direct effect on best fitness is only
$\eta^2 \approx 0.02$. Topology acts primarily through the
diversity channel, not by directly improving or degrading the best
solution found.

\paragraph{Interpretation.}
Simple tasks are structurally smooth---all paths lead to the same
optimum, so topology barely matters. Complex tasks are
rugged---multiple competing strategies exist, and topology determines
whether the system explores many of them (sparse connectivity) or
collapses to one (dense connectivity). Real-world agent deployments
overwhelmingly involve the latter: coding tasks with multiple valid
algorithms, planning tasks with multiple valid orderings, analysis
tasks with multiple valid framings. The $14\times$ amplification from
$K = 0$ to $K = 6$ suggests that topology effects in production
systems are far larger than simple benchmarks would indicate.

%% ===================================================================
\section{Experiment 3: Ring vs Star at Constant $\beta_1$}
\label{sec:exp3}
%% ===================================================================

Experiments~1--2 vary cycle count across topologies with different
structures. But a natural question arises: within a single cycle rank
class ($\beta_1 = 1$), how much does the \emph{arrangement} of edges
matter? A 5-node ring and a 5-node star both have $\beta_1 = 1$ and
$|E| = 5$, but very different degree distributions: the ring is
2-regular while the star has one hub (degree~4) and four leaves
(degree~1). If cycle rank is the primary driver, these should behave
similarly despite their structural differences.

\paragraph{Design.}
Two topologies on $n = 5$ nodes: ring (2-regular, diameter~2,
$\lambda_2 = 0.382$) and star (irregular, diameter~2,
$\lambda_2 = 1.0$). Both have $\beta_1 = 1$ and $|E| = 5$. NK
landscapes with $K \in \{0, 2, 4, 6\}$, 20~seeds per condition,
160~runs total (80~ring, 80~star). Island-model GA: 50 individuals
per island, 500~generations.

\paragraph{Results.}
The ring-vs-star effect is small compared to the cross-$\beta_1$
effects in Experiments~1--2 (Table~\ref{tab:ring-star}). At
generation~500:

\begin{table}[t]
\centering
\caption{Ring vs star comparison at constant $\beta_1 = 1$. Effect
sizes are small ($\eta^2 < 0.15$) and mostly non-significant,
contrasting with $\eta^2 = 0.17$--$0.69$ across $\beta_1$ classes in
Experiments~1--2. Cycle rank matters more than degree distribution.}
\label{tab:ring-star}
\begin{tabular}{@{}cccccl@{}}
\toprule
$K$ & $\eta^2$ (diversity) & $p$ & Ring mean & Star mean & Winner \\
\midrule
0 & 0.014 & 0.59 & 0.046 & 0.047 & --- \\
2 & 0.132 & 0.22 & 0.062 & 0.045 & ring \\
4 & 0.077 & 0.10 & 0.148 & 0.103 & ring \\
6 & 0.080 & 0.12 & 0.249 & 0.192 & ring \\
\bottomrule
\end{tabular}
\end{table}

% FIGURE: ring-vs-star-diversity.pdf
% Caption: Ring vs star diversity at generation 500 across NK landscape
% ruggedness. Both topologies have beta_1 = 1. The ring consistently
% maintains slightly higher diversity, but the effect is small (eta^2
% = 0.08--0.13) compared to cross-beta_1 effects.

\paragraph{Temporal dynamics.}
The ring-vs-star effect is most pronounced at intermediate
generations. At $K = 6$, generation~100: $\eta^2 = 0.161$
($p = 0.02$, the only significant comparison). By generation~500,
this fades to $\eta^2 = 0.080$ ($p = 0.12$). The ring's lower
$\lambda_2$ (0.382 vs 1.0) slows information propagation, preserving
diversity transiently but not permanently.

\paragraph{Interpretation.}
The ring consistently maintains slightly higher diversity than the
star, favoring ring in all $K > 0$ conditions. The mechanism is
$\lambda_2$: the star's higher algebraic connectivity ($\lambda_2 =
1.0$ vs $0.382$) enables faster consensus, collapsing diversity more
rapidly. But the effect is modest ($\eta^2 \leq 0.13$, none
significant at $\alpha = 0.05$ after generation~100) compared to the
large effects seen when varying $\beta_1$ ($\eta^2 = 0.17$--$0.69$).
This confirms that \textbf{cycle rank is the primary topology
invariant for diversity control}, with degree distribution and
$\lambda_2$ playing a secondary role.

The plateau pattern is also informative: the ring advantage does
\emph{not} grow monotonically with $K$, instead leveling off at
$\eta^2 \approx 0.08$--$0.13$ for $K \geq 2$. Once the landscape is
rugged enough to create competing strategies, the within-$\beta_1$
topology effect stabilizes. This contrasts sharply with the
cross-$\beta_1$ NK dose-response (Experiment~2), where $\eta^2$ grows
from 0.05 to 0.69.

%% ===================================================================
\section{Experiment 4: Bridge --- $\beta_1$ vs $\lambda_2$ Decomposition}
\label{sec:exp4}
%% ===================================================================

Experiments~1--3 establish that cycle rank $\beta_1$ is the primary
invariant, with Experiment~3 hinting that $\lambda_2$ plays a
secondary role. But $\beta_1$ and $\lambda_2$ are themselves
correlated: adding edges to a graph typically increases both. Can we
disentangle them?

We construct \emph{iso-spectral graph families}: sets of graphs that
share the same $\lambda_2$ but differ in $\beta_1$. Within a family,
any diversity difference must be driven by $\beta_1$ (since
$\lambda_2$ is held constant). \emph{Across} families (different
$\lambda_2$, same $\beta_1$), any difference must be driven by
$\lambda_2$.

\paragraph{Design.}
Two families of four graphs each, all on $n = 5$ vertices:

\begin{itemize}[nosep]
  \item \textbf{Cycle family} ($\lambda_2 \approx 0.586$): 5-node
    ring ($\beta_1 = 1$), ring + 1~chord ($\beta_1 = 2$), ring +
    2~chords ($\beta_1 = 3$), ring + 3~chords ($\beta_1 = 4$). Adding
    chords increases $\beta_1$ while preserving the spectral gap via
    careful edge selection.

  \item \textbf{Star family} ($\lambda_2 = 1.0$): star with pendant
    ($\beta_1 = 1$), star + 1~leaf-edge ($\beta_1 = 2$), star +
    2~leaf-edges ($\beta_1 = 3$), star + 3~leaf-edges ($\beta_1 = 4$).
    Adding leaf-to-leaf connections increases $\beta_1$ while the
    hub structure preserves $\lambda_2$.
\end{itemize}

\noindent
Two NK landscapes: $K = 0$ (smooth, negative control) and $K = 4$
(rugged). 20~seeds per condition. Total: $8 \text{ topologies} \times
2 \text{ landscapes} \times 20 \text{ seeds} = 320$~runs.

\paragraph{Results.}
The experiment reveals two independent mechanisms operating at
different timescales (Table~\ref{tab:bridge}):

\begin{table}[t]
\centering
\caption{Bridge experiment: two-way ANOVA decomposing diversity
variance into $\beta_1$ (within-family) and $\lambda_2$
(cross-family) effects on NK $K = 4$ landscapes. $\beta_1$ peaks
early and fades; $\lambda_2$ emerges later and persists.}
\label{tab:bridge}
\begin{tabular}{@{}ccccc@{}}
\toprule
Generation & $\eta^2_{\beta_1}$ & $p_{\beta_1}$ &
$\eta^2_{\lambda_2}$ & $p_{\lambda_2}$ \\
\midrule
50  & 0.037 & 0.11 & 0.001 & 0.76 \\
100 & 0.057 & 0.03 & 0.006 & 0.41 \\
200 & 0.045 & 0.07 & 0.112 & $< 0.001$ \\
500 & 0.014 & 0.44 & 0.066 & 0.005 \\
\bottomrule
\end{tabular}
\end{table}

% FIGURE: bridge-diversity-over-generations.pdf
% Caption: Bridge experiment: diversity trajectories for two iso-spectral
% families. Within each family (constant lambda_2, varying beta_1), higher
% beta_1 produces lower diversity at early generations (transient effect).
% Across families (different lambda_2, same beta_1), the star family
% (higher lambda_2) converges faster (persistent effect). Two invariants,
% two timescales.

\paragraph{$\beta_1$ drives transient diversity.}
Within the cycle family ($\lambda_2 = 0.586$), $\beta_1$ explains
$\eta^2 = 0.191$ of diversity variance at generation~100
($F(3, 76) = 5.997$, $p = 0.001$) but fades to $\eta^2 = 0.067$ by
generation~500 ($p = 0.15$). Higher $\beta_1$ means more feedback
loops, which accelerate the initial mixing of information---but this
advantage is transient. The star family ($\lambda_2 = 1.0$) shows a
similar but weaker within-family $\beta_1$ effect ($\eta^2 = 0.092$
at generation~100, $p = 0.06$).

\paragraph{$\lambda_2$ drives persistent diversity.}
Across families at matched $\beta_1$, the family effect
($\lambda_2$) is negligible at generation~50 ($\eta^2 = 0.001$) but
emerges strongly at generation~200 ($\eta^2 = 0.112$,
$p < 0.001$) and persists through generation~500 ($\eta^2 = 0.066$,
$p = 0.005$). The cross-family analysis at $\beta_1 = 1$ confirms:
$\eta^2 = 0.265$ at generation~200 ($p = 0.003$), meaning the star
family ($\lambda_2 = 1.0$) converges to lower diversity than the
cycle family ($\lambda_2 = 0.586$) at the same cycle rank.

\paragraph{Negative control.}
On the smooth landscape ($K = 0$), both effects are negligible:
$\eta^2_{\beta_1} < 0.02$ and $\eta^2_{\lambda_2} < 0.02$ at all
generations. This confirms that both mechanisms require landscape
ruggedness to produce measurable effects, consistent with the
dose-response in Experiment~2.

\paragraph{Interpretation.}
$\beta_1$ and $\lambda_2$ are not redundant---they control
diversity through two independent mechanisms on two timescales.
$\beta_1$ counts feedback loops, which determine how quickly initial
information is mixed (transient effect, first $\sim$100 generations).
$\lambda_2$ governs the spectral gap, which determines how quickly
the system reaches global consensus (persistent effect, visible from
generation~200 onward). For agent systems, this suggests:
\emph{cycle rank controls the exploration window} (how long the
system entertains diverse strategies), while \emph{algebraic
connectivity controls the convergence rate} (how quickly the system
settles). These are independent design levers.

%% ===================================================================
\section{Experiment 5: LLM Sign-Flip (Claude Chorus)}
\label{sec:exp5}
%% ===================================================================

The GA experiments (Sections~\ref{sec:exp1}--\ref{sec:exp4}) establish
that topology controls diversity through cycle rank and spectral gap.
But does this transfer to LLM agents? The node-level differences
(populations vs.\ single outputs, genetic operators vs.\ language
generation) are substantial. We test directly.

\paragraph{Design.}
Five GPT-4o-mini agents, eight coding tasks (shortest path, kth
largest, rate limiter, JSON parser, LRU cache, regex matcher,
expression evaluator, Bloom filter), four topologies (none, ring,
star, fully connected), 10~seeds per condition. Two coupling
conditions, 1,600 API calls each (3,200 total):

\begin{itemize}[nosep]
  \item \textbf{Negative coupling:} each agent is instructed to use a
    \emph{different} algorithm from those visible in its context.
  \item \textbf{Positive coupling:} each agent is instructed to
    ``build on and improve'' neighbors' solutions.
\end{itemize}

\noindent
Diversity is measured as mean pairwise cosine distance on sentence
embeddings (all-MiniLM-L6-v2) of agent outputs.

\paragraph{Metric validation.}
The ``none'' topology provides a sanity check: all five agents
generate independently with no information sharing, yet produce
near-identical outputs (mean diversity 0.090, $\sigma \approx 0.01$).
This confirms the metric detects the known homogeneity of LLM prior
distributions---LLMs start from a peaked prior, unlike GAs' random
initialization. The low within-condition variance also explains the
large Cohen's $d$ values below.

\paragraph{Results: negative coupling.}
The diversity ordering \textbf{inverts} from the GA result
(Table~\ref{tab:signflip}):

\begin{table}[t]
\centering
\caption{Topology--diversity ordering inverts between GAs and LLM
agents. The sign depends on whether coupling merges solutions
(GA-like) or constrains them (``be different''). A single instruction
change flips the ordering.}
\label{tab:signflip}
\begin{tabular}{@{}lll@{}}
\toprule
\textbf{System} & \textbf{Diversity ordering} & \textbf{Statistic} \\
\midrule
GA & none $>$ ring $>$ FC & $W = 1.0$ (6 domains) \\
LLM (negative) & FC $>$ ring $\approx$ star $>$ none &
  $W = 0.91$ (6/8 tasks) \\
LLM (positive) & none $>$ star $\approx$ ring $>$ FC &
  $\eta^2 = 0.291$ \\
\bottomrule
\end{tabular}
\end{table}

Six of eight tasks show the inverted ordering cleanly (Kruskal-Wallis
$p < 10^{-6}$; Cohen's $d$ from $-14.0$ to $-1.9$). Two exceptions
are informative:

\begin{itemize}[nosep]
  \item \textbf{Bloom filter} ($d \approx 0$): a unimodal task with
    essentially one canonical implementation. No algorithmic variety
    for topology to act on.
  \item \textbf{Expression evaluator} ($d = +0.4$, positive): the
    model's prior is already spread across multiple parsing strategies
    (recursive descent, stack-based, Pratt). When the prior is
    already flat, isolation produces diverse solutions---the GA
    regime---and negative coupling adds little.
\end{itemize}

% FIGURE: llm-diversity-heatmap.pdf
% Caption: LLM diversity across 8 coding tasks and 4 topologies under
% negative coupling (left) and positive coupling (right). Under negative
% coupling, connectivity increases diversity (FC highest). Under positive
% coupling, connectivity decreases diversity (none highest). The sign
% flip is clean across 6/8 tasks.

\paragraph{Why the sign flips.}
Under ``none'' topology, all five LLM agents independently generate
the \emph{same} algorithm (e.g., quickselect for kth\_largest). LLMs
start from a \textbf{peaked prior} (low entropy); GAs start from
\textbf{random initialization} (high entropy). Isolation preserves
initial entropy: high for GAs (diverse), low for LLMs (homogeneous).

\begin{itemize}[nosep]
  \item \textbf{GA migration} is colimit-like: it merges populations
    by injecting solutions into neighboring islands. More connectivity
    $\to$ faster homogenization $\to$ less diversity.
  \item \textbf{LLM ``be different'' context} is limit-like: it
    accumulates exclusion constraints from neighbors. More connectivity
    $\to$ more constraints $\to$ agents forced into distributional
    tails $\to$ \emph{more} diversity.
\end{itemize}

\noindent
The topology is the invariant; the sign depends on the coupling
direction. We conjecture this sign reversal is an instance of the
adjunction between colimit (merge) and limit (constrain) along the
topology functor, with the formal categorical construction developed
in a companion paper~\cite{langer2026games}.

\paragraph{Positive coupling confirms the prediction.}
With positive coupling (``build on and improve''), 1,600 additional
API calls. Results: the diversity ordering \textbf{reverts to
GA-like}: none $>$ star $\approx$ ring $>$ FC (Kruskal-Wallis
$H = 94.83$, $p = 2.01 \times 10^{-20}$; Cohen's $d = 1.08$,
none vs.\ FC; $\eta^2 = 0.291$). The ``none'' topology shows
identical diversity across coupling conditions (negative: 0.090,
positive: 0.095, $p = 0.96$), confirming that coupling direction, not
some other variable, drives the reversal. Connected topologies
collapse to near-zero diversity (0.025--0.028) under positive
coupling. A single instruction change flips the diversity ordering.

%% ===================================================================
\section{Discussion}
\label{sec:discussion}
%% ===================================================================

\subsection{The Evidence Gap}

At least nine recent agent systems---spanning ICML, AAAI, WWW,
NeurIPS, and AAMAS---chose topologies without density-controlled
comparisons (Table~\ref{tab:evidence-gap}). Our confound theorem
(Theorem~\ref{thm:confound}) shows these comparisons are
fundamentally uninterpretable: they cannot distinguish density effects
from cycle effects.

\begin{table}[t]
\centering
\caption{Agent systems whose topology choices lack controlled
evidence. ``Confound type'' indicates which properties are varied
simultaneously: D = density, C = cycle structure, S = spectral gap.}
\label{tab:evidence-gap}
\footnotesize
\begin{tabular}{@{}llp{4.5cm}@{}}
\toprule
\textbf{System} & \textbf{Confound} & \textbf{Issue} \\
\midrule
Graph-GRPO~\cite{graphgrpo2026} & D+C & Complete graph = worst
  performance; also max $\beta_1$ \\
OFA-MAS~\cite{ofamas2026} & D+C & Prunes density \& cycles together;
  generates only DAGs \\
G-Designer~\cite{gdesigner2025} & D+C+S & No density-controlled
  baseline \\
GoAgent~\cite{goagent2026} & D+C & DAG vs.\ dense alternatives
  confounded \\
ARG-Designer~\cite{argdesigner2026} & C & DAG search space only;
  cycles never considered \\
AgentConductor~\cite{agentconductor2026} & C & Cycles never tested \\
DyTopo~\cite{dytopo2026} & C & DAG enforced throughout; dynamic
  topology shifts confound time \& structure \\
HERA~\cite{hera2026} & D & Cycles tracked but not controlled;
  density varies with cycle count \\
Dochkina~\cite{dochkina2026} & D+C+S & 25K tasks, $d = 1.86$, but
  protocols differ on all dimensions \\
\bottomrule
\end{tabular}
\end{table}

The DAG hegemony is striking: of these nine systems, six are
explicitly restricted to DAG topologies ($\beta_1 = 0$), and two more
implicitly favor sparse topologies. Only HERA~\cite{hera2026} tracks
cycle count as a design metric, and even it does not control density.
The implicit assumption---that acyclic structures are
preferable---has never been tested against the alternative.

\subsection{The Invariant Hierarchy}

Our five experiments reveal three graph invariants operating at
different timescales:

\begin{enumerate}[nosep]
  \item \textbf{Directed cycle count $\kappa$} (Experiment~1):
    determines diversity at constant density. Effect visible from
    early generations. Captures the \emph{arrangement} of feedback.

  \item \textbf{Cycle rank $\beta_1$} (Experiments~2--4): the primary
    invariant for diversity control. Drives transient diversity
    ($\eta^2 = 0.191$ at generation~100, fading by generation~500 in
    Experiment~4). Captures the \emph{amount} of feedback.

  \item \textbf{Algebraic connectivity $\lambda_2$} (Experiments~3--4):
    drives persistent diversity ($\eta^2 = 0.112$--$0.265$ at
    generation~200+, not fading). Captures the \emph{rate} of
    consensus. Within a single $\beta_1$ class (Experiment~3), ring
    ($\lambda_2 = 0.382$) maintains more diversity than star
    ($\lambda_2 = 1.0$).
\end{enumerate}

% FIGURE: invariant-hierarchy.pdf
% Caption: The three-invariant hierarchy. Directed cycle count kappa
% determines the arrangement of feedback (fast). Cycle rank beta_1
% determines the amount of feedback, driving transient diversity
% (medium). Algebraic connectivity lambda_2 determines consensus rate,
% driving persistent diversity (slow). Each operates at a distinct
% timescale.

\noindent
Whether these form a nested hierarchy ($\kappa$ modulates $\beta_1$
which modulates $\lambda_2$) or operate independently on distinct
phenotypic properties is an open question. The bridge experiment
(Experiment~4) shows that $\beta_1$ and $\lambda_2$ produce
\emph{independent} effects (low cross-correlation in the two-way
ANOVA), suggesting the latter interpretation. Future work should test
whether ablating one invariant collapses the other's effect (nested)
or degrades a different metric independently (parallel).

\subsection{Coupling Sign and Agent Design}

The LLM sign-flip (Experiment~5) has direct implications for agent
orchestration:

\begin{itemize}[nosep]
  \item \textbf{Know your coupling direction.} A merging coupling
    (``build on previous answers'') homogenizes like GA migration:
    more connectivity $\to$ less diversity. A constraining coupling
    (``do something different'') diversifies: more connectivity $\to$
    more diversity. The same star topology can either preserve or
    destroy diversity depending on a single instruction.

  \item \textbf{LLM priors are peaked.} Independent agents (no
    topology) produce near-identical outputs. ``Ensemble diversity''
    claims based on running the same model independently are illusory
    when models share a peaked prior.

  \item \textbf{Error amplification is structural.} Moran's
    $17\times$ error amplification~\cite{moran2025error} in
    unstructured communication arises from high $\beta_1$ (fully
    connected) combined with merging coupling. The standard fix (impose
    DAG structure) simultaneously removes cycles \emph{and} density---
    we cannot know which helps.
\end{itemize}

\subsection{Practical Recommendations}

\begin{enumerate}
  \item \textbf{Treat topology as a first-class design variable.}
    Topology controls the diversity--convergence tradeoff in both GAs
    and LLM agents ($\eta^2$ up to 0.69 in GAs, 0.291 in LLMs).
    Dochkina's 44\% quality spread across 25K tasks confirms the
    effect at industrial scale.

  \item \textbf{Know your coupling direction.} A single instruction
    change flips which topology produces the most diverse output.
    Merging homogenizes; differentiating diversifies.

  \item \textbf{Match topology to task complexity.} Simple tasks
    tolerate any topology ($\eta^2 = 0.05$ at $K = 0$); complex
    tasks demand deliberate choice ($\eta^2 = 0.69$ at $K = 6$).

  \item \textbf{Use $\beta_1$ for exploration budget, $\lambda_2$
    for convergence rate.} These are independent levers. More cycles
    extend the exploration window; lower $\lambda_2$ slows consensus.

  \item \textbf{Monitor diversity, not just accuracy.} Topology
    effects manifest in strategy diversity. Accuracy alone masks
    topology-induced failures.
\end{enumerate}

\subsection{Limitations}

Our GA experiments control density precisely using iso-dense directed
graphs (Experiment~1) and iso-spectral undirected families
(Experiment~4). Our LLM experiment (Experiment~5) does \emph{not}
control density: the four topologies (none, ring, star, fully
connected) differ in both density and structure. The \emph{sign-flip}
finding is robust to this confound---both coupling conditions use
identical topologies, and the reversal is within-topology---but the
absolute diversity ordering across LLM topologies is confounded by
density. Density-controlled LLM experiments using iso-dense directed
topologies are needed.

Additionally, our diversity metric (cosine distance on sentence
embeddings) is validated only implicitly via the ``none'' baseline. A
systematic comparison against human-judged algorithmic similarity
would strengthen the measurement. Our GA experiments use $N = 20$ bit
strings; scaling to realistic problem sizes is an open question. The
bridge experiment iso-spectral families are only approximately
matched ($\lambda_2 = 0.586$ vs $1.0$); exact spectral matching
across families with different $\beta_1$ ranges requires construction
techniques we leave to future work.

%% ===================================================================
\section{Related Work}
\label{sec:related}
%% ===================================================================

\paragraph{Evolutionary computation.}
The topology--diversity link in island-model GAs is well
established~\cite{tomassini2005spatially,skolicki2005influence}, but
prior work does not control for density. Tomassini's comprehensive
treatment~\cite{tomassini2005spatially} analyzes grid, ring, and
random topologies without iso-dense controls. Our Experiment~1 is, to
our knowledge, the first density-controlled topology comparison in
evolutionary computation.

\paragraph{LLM multi-agent systems.}
The field is converging on topology as a key design
variable~\cite{yang2025topology}, but comparisons remain
uncontrolled. Graph-GRPO~\cite{graphgrpo2026} learns topology via
reinforcement learning, finding that complete graphs perform worst---
consistent with our $\beta_1$ prediction, but confounded by density.
OFA-MAS~\cite{ofamas2026} generates topologies via mixture-of-experts
models but restricts to DAGs. G-Designer~\cite{gdesigner2025} uses
graph neural networks to architect topologies without
density-controlled baselines. GoAgent~\cite{goagent2026} introduces
group-level topology but compares DAGs against dense alternatives.
ARG-Designer~\cite{argdesigner2026} and
AgentConductor~\cite{agentconductor2026} operate exclusively in DAG
space. DyTopo~\cite{dytopo2026} dynamically adjusts topology during
execution---an important direction, but enforces DAG structure
throughout.

HERA~\cite{hera2026} is the closest to our approach: it tracks cycle
count during topology evolution, finding that ``compact chains with
strategically preserved cycles'' achieve 38.69\% average improvement.
This is the first system to treat cycle count as a topology metric,
but it does not control density and cannot distinguish cycle effects
from density effects.

Dochkina~\cite{dochkina2026} provides the largest-scale topology
comparison to date: 25,000 tasks, 256 agents, 8 coordination
protocols, with a 44\% quality spread ($d = 1.86$). Crucially, the
effect is capability-moderated: weak models prefer DAG-like
hierarchies, while strong models can exploit cyclic coordination
structures. This supports our finding that topology effects are
amplified on hard problems (Experiment~2): stronger models on harder
tasks correspond to higher effective $K$.

Wang et al.~\cite{wang2026kregular} show that $K$-regular graphs
maximize coordination efficiency, with a minimum capability threshold
$\beta_{\min} \sim 1/K$. This suggests a \emph{coordination--exploration
duality}: density helps coordination but hurts exploration (our
finding). The optimal topology balances these competing demands---a
Goldilocks $\beta_1$.

\paragraph{Categorical and algebraic approaches.}
Capucci and Myers~\cite{capucci2026feedback} formalize agent systems
as lenses and Moore machines, showing that each cycle in the
communication graph requires an independent convergence guarantee.
This provides the theoretical cost side of our empirical benefit
measurements: cycles buy exploration (our Experiments~1--4) but cost
convergence guarantees. Langer and Turing~\cite{langer2026games}
develop the categorical framework (lax monoidal functors, Kleisli
morphisms over an evolution monad) that motivates the colimit/limit
interpretation of the coupling-sign flip; our Experiment~5 provides
the first empirical test.

%% ===================================================================
\section{Conclusion}
\label{sec:conclusion}
%% ===================================================================

Communication topology is not neutral infrastructure---it is a design
variable controlling how multi-agent systems balance exploration
against convergence. We have shown:

\begin{enumerate}[nosep]
  \item \textbf{The confound is real.} Cycle rank and density are
    algebraically identical for undirected graphs. Most published
    topology comparisons are uninterpretable.

  \item \textbf{Topology independently controls diversity.}
    $\eta^2 = 0.17$ at constant density (Experiment~1), scaling to
    0.69 on hard problems (Experiment~2).

  \item \textbf{Two invariants, two timescales.} $\beta_1$ drives
    transient exploration; $\lambda_2$ drives persistent convergence.
    These are independent levers (Experiment~4).

  \item \textbf{Cycle rank $>$ degree distribution.} Within a single
    $\beta_1$ class, degree differences explain at most
    $\eta^2 \approx 0.13$ (Experiment~3).

  \item \textbf{The sign depends on coupling.} A single instruction
    change flips which topology produces diverse output
    ($\eta^2 = 0.291$, Experiment~5). GA intuition has the wrong sign
    for constraining agents.
\end{enumerate}

\noindent
For safe multi-agent deployment, topology must be selected
deliberately, its coupling direction understood, and its effects
monitored through diversity metrics. The three-invariant hierarchy
($\kappa \to \beta_1 \to \lambda_2$) provides a principled framework
for this selection. Current systems are tuning their most powerful
design variable blindly.

\smallskip
\noindent\textbf{Acknowledgments.} [Withheld for review.]

%% ===================================================================
%% References
%% ===================================================================

\bibliographystyle{splncs04}
\begin{thebibliography}{20}

\bibitem{capucci2026feedback}
Capucci, M., Myers, D.J.:
Feedback in categorical multi-agent systems.
arXiv:2604.03017 (2026)

\bibitem{dochkina2026}
Dochkina, V.:
Drop the hierarchy and roles: How self-organizing {LLM} agents
outperform designed structures.
arXiv:2603.28990 (2026)

\bibitem{dytopo2026}
Lu, Y., et al.:
DyTopo: Dynamic topology routing for multi-agent reasoning via
semantic matching.
arXiv:2602.06039 (2026)

\bibitem{fiedler1973algebraic}
Fiedler, M.:
Algebraic connectivity of graphs.
Czechoslovak Mathematical Journal \textbf{23}(2), 298--305 (1973)

\bibitem{goagent2026}
Chen, H., et al.:
GoAgent: Group-of-agents communication topology generation for
{LLM}-based multi-agent systems.
arXiv:2603.19677 (2026)

\bibitem{gdesigner2025}
Zhang, G., et al.:
G-Designer: Architecting multi-agent communication topologies via
graph neural networks.
In: Proc.\ ICML (2025)

\bibitem{graphgrpo2026}
Cang, Y., et al.:
Graph-GRPO: Stabilizing multi-agent topology learning via group
relative policy optimization.
arXiv:2603.02701 (2026)

\bibitem{hatcher2002algebraic}
Hatcher, A.:
Algebraic Topology.
Cambridge University Press (2002)

\bibitem{hera2026}
Li, S., Ramakrishnan, N.:
Experience as a compass: Multi-agent {RAG} with evolving
orchestration and agent prompts.
arXiv:2604.00901 (2026)

\bibitem{johnson1975finding}
Johnson, D.B.:
Finding all the elementary circuits of a directed graph.
SIAM J.\ Comput.\ \textbf{4}(1), 77--84 (1975)

\bibitem{kauffman1993origins}
Kauffman, S.A.:
The Origins of Order.
Oxford University Press (1993)

\bibitem{langer2026games}
Langer, R., Turing, C.:
From games to graphs: Categorical composition of genetic algorithms
across domains.
In: Proc.\ ACT (2026)

\bibitem{agentconductor2026}
Wang, S., et al.:
AgentConductor: Topology evolution for multi-agent competition-level
code generation.
arXiv:2602.17100 (2026)

\bibitem{argdesigner2026}
Li, S., et al.:
Assemble your crew: Automatic multi-agent communication topology
design via autoregressive graph generation.
In: Proc.\ AAAI (2026)

\bibitem{moran2025error}
Moran, K.:
The secret is topology: Why unstructured multi-agent communication
amplifies errors.
Towards Data Science (2025)

\bibitem{ofamas2026}
Li, S., et al.:
OFA-MAS: One-for-all multi-agent system topology design based on
mixture-of-experts graph generative models.
In: Proc.\ WWW (2026)

\bibitem{skolicki2005influence}
Skolicki, Z., De~Jong, K.A.:
The influence of migration sizes and intervals on island models.
In: GECCO, pp.\ 1295--1302 (2005)

\bibitem{tomassini2005spatially}
Tomassini, M.:
Spatially Structured Evolutionary Algorithms.
Springer (2005)

\bibitem{wang2026kregular}
Wang, X., et al.:
$K$-regular graphs maximize coordination in multi-agent systems.
arXiv:2604.07751 (2026)

\bibitem{yang2025topology}
Yang, J., et al.:
Topological structure learning should be a research priority for
{LLM}-based multi-agent systems.
arXiv:2505.22467 (2025)

\end{thebibliography}

\end{document}
